\chapter{Objectives}

\section[Primary Objectives]{\textbf{Primary Objectives}}

\subsection{Main Objectives}

The main objectives of the EdgeCare Response platform are:
\begin{enumerate}
\item To design and implement a low-power wearable edge node (ESP32, MPU6050, heart-rate sensor) capable of continuous inertial and physiological monitoring with on-device TinyML fall prediction, targeting $\geq$95\% accuracy, recall, and specificity under subject-wise testing.
\item To develop a hybrid Edge-AI fall verification engine combining quantized TinyML inference with deterministic physical thresholds and personalized physiological heart-rate deviation analysis to reduce false alarms compared to model-only or threshold-only baselines.
\item To design an Emergency Severity Score (ESS) model that integrates fall probability, peak acceleration, angular velocity, heart-rate deviation, inactivity duration, and user response status to classify events as Mild, Moderate, or Critical for tiered emergency escalation.
\item To build an integrated emergency coordination workflow connecting caregiver notification, ambulance dispatch with real-time GPS tracking, and hospital pre-notification — delivered through a cloud-based Next.js dashboard with role-based access for caregivers, dispatchers, and hospital staff.
\end{enumerate}

\subsection{Expected Outcomes}

The expected outcomes from achieving the main objectives are:
\begin{itemize}
\item A fully functional wearable prototype achieving $\geq$95\% fall detection accuracy with $<$100 ms on-device inference per window and 24--72 hour battery life under adaptive duty cycling.
\item A hybrid verification engine with significantly reduced false alarm escalation rates compared to model-only and threshold-only baseline systems.
\item An ESS triage model correctly classifying Mild, Moderate, and Critical events, enabling context-appropriate automated responses without unnecessary ambulance dispatch for low-risk events.
\item A real-time cloud dashboard enabling caregivers to receive alerts in $<$10 seconds and critical dispatch coordination to be initiated in $<$30 seconds after fall verification.
\end{itemize}

\section[Secondary Objectives]{\textbf{Secondary Objectives}}

\subsection{Additional Objectives}

\begin{enumerate}
\item To implement an adaptive duty-cycling strategy on the ESP32 that balances monitoring accuracy with energy efficiency, extending battery life through coordinated sleep/sample/infer/transmit state management.
\item To validate the system's end-to-end performance using public inertial datasets (SisFall, MobiAct, KFall) and simulated emergency workflow scenarios across caregiver, dispatch, ambulance, and hospital states.
\end{enumerate}

\subsection{Supporting Goals}

The following supporting goals complement the primary and secondary objectives:
\begin{itemize}
\item Hands-on experience with embedded systems programming (ESP32, Arduino/ESP-IDF) and I2C sensor interfacing (MPU6050, heart-rate sensor).
\item Practical knowledge of TinyML model development, training, 8-bit post-training quantization, and on-device deployment using TensorFlow Lite Micro within strict memory and latency constraints.
\item Understanding of event-driven cloud microservices architecture using Python FastAPI, Redis Pub/Sub, and PostgreSQL for healthcare data management.
\item Proficiency in real-time frontend development using Next.js, TypeScript, WebSockets, and Leaflet.js for live ambulance GPS tracking.
\item Experience with role-based access control (RBAC) and secure session management (Google OAuth 2.0, JWT HttpOnly cookies) for multi-stakeholder healthcare platforms.
\item Knowledge of external API integration including Google Maps Directions API (routing and ETA), Twilio (SMS/voice alerts), and Brevo SMTP (email notifications).
\item Understanding of physiological signal processing concepts including personalized heart-rate baseline computation, normalized deviation scoring, and inactivity detection.
\item Awareness of privacy, consent, and ethical governance requirements for wearable health monitoring systems intended for older adult populations.
\item Collaborative embedded-to-cloud development practices including hardware-software co-design, iterative prototyping, and cross-functional team coordination.
\item Foundation for future ethics-approved clinical pilot studies, longitudinal gait analytics, and EMS-integrated smart ambulance allocation.
\end{itemize}